\chapter{Implementation}
\label{chap:implementation}

The implementation followed the project evolution described in Chapter~2. It
started from the company crawler baseline, used n8n to validate a browser-agent
fallback, and then formalized the maintainability direction through OWC. This
chapter describes what was implemented without claiming that the full
Python/FastAPI/Puppeteer platform replaced the existing company crawler.

\section{Existing Crawler Baseline}

\subsection{Company Puppeteer Crawlers}

The company already used Puppeteer-based browser automation to inspect known
streaming website patterns. Those crawlers could open targets, follow expected
paths, and collect useful evidence when the page structure matched previous
rules.

\subsection{What the Existing Crawlers Already Solved}

The baseline solved the operational need for repeatable collection on stable
patterns. It provided a practical starting point because the company already
understood browser automation, page traversal, and evidence handoff.

\subsection{Maintenance and Adaptation Limits}

The maintenance burden appeared when websites changed selectors, popups,
iframes, server controls, redirect paths, or stream-generation logic. A fixed
rule could fail even when a human reviewer could still recover the evidence by
observing the page and trying a different action.

\subsection{Why a Browser-Agent Layer Was Explored}

A browser-agent layer was explored to make the fallback path more adaptable and
traceable. The goal was not to remove crawler engineering, but to support it
when a known pattern stopped working and the page still contained recoverable
evidence.

\section{Prototype Phase with n8n}

\subsection{Prototype Goal}

The first objective was speed. n8n was used to validate whether an LLM-guided
browser workflow could classify streaming pages, navigate through landing and
hosting pages, activate players, and return useful evidence before investing in
a stricter architecture~\cite{n8n}.

\subsection{Why n8n Was Used}

n8n made it easy to connect browser actions, agent steps, intermediate outputs,
and manual inspection. That made it useful for experimentation even though it
was not the final architecture.

\subsection{n8n Landing Agent}

In the company context, the n8n landing agent was deployed as a backup solution
to support the existing crawler workflow. Its role was not to replace the full
crawler system. It provided a flexible fallback when crawler rules failed to
recover useful landing-page candidates after website changes.

\begin{figure}[!htbp]
\centering
\begingroup
\setlength{\fboxsep}{3pt}
\setlength{\fboxrule}{0.35pt}
\fcolorbox{gray!35}{white}{\includegraphics[width=0.92\textwidth]{\detokenize{n8n Full Workflow.png}}}
\endgroup

{\scriptsize\bfseries Full n8n workflow}

\vspace{0.45em}
\begin{minipage}{0.48\textwidth}
  \centering
  \begingroup\setlength{\fboxsep}{2pt}\setlength{\fboxrule}{0.3pt}%
  \fcolorbox{gray!35}{white}{\includegraphics[width=0.94\linewidth]{\detokenize{n8n classification agent.png}}}\endgroup\\[-0.1em]
  {\scriptsize\bfseries Classification agent}
\end{minipage}\hfill
\begin{minipage}{0.48\textwidth}
  \centering
  \begingroup\setlength{\fboxsep}{2pt}\setlength{\fboxrule}{0.3pt}%
  \fcolorbox{gray!35}{white}{\includegraphics[width=0.94\linewidth]{\detokenize{n8n landing page.png}}}\endgroup\\[-0.1em]
  {\scriptsize\bfseries Landing-page agent}
\end{minipage}

\vspace{0.55em}
\begin{minipage}{0.48\textwidth}
  \centering
  \begingroup\setlength{\fboxsep}{2pt}\setlength{\fboxrule}{0.3pt}%
  \fcolorbox{gray!35}{white}{\includegraphics[width=0.94\linewidth]{\detokenize{n8n Hosting page.png}}}\endgroup\\[-0.1em]
  {\scriptsize\bfseries Hosting-page agent}
\end{minipage}\hfill
\begin{minipage}{0.48\textwidth}
  \centering
  \begingroup\setlength{\fboxsep}{2pt}\setlength{\fboxrule}{0.3pt}%
  \fcolorbox{gray!35}{white}{\includegraphics[width=0.94\linewidth]{\detokenize{n8n Embedded Agent.png}}}\endgroup\\[-0.1em]
  {\scriptsize\bfseries Embedded-player agent}
\end{minipage}

\caption{n8n workflow screenshots: full sequential workflow and the four specialist agent workflows.}
\label{fig:n8n-agent-screenshots}
\end{figure}


\subsection{Validation on Real Websites}

The prototype was tested on real streaming pages with landing listings,
hosting pages, embedded players, popups, server switches, and delayed media
requests. These tests confirmed that a browser agent could recover evidence in
some cases, but also showed that long browser workflows need stronger state and
traceability.

\subsection{Prototype Findings}

\begin{table}[!htbp]
\centering
\begin{tabularx}{\textwidth}{Q{3.1cm} Y Y}
\toprule
\rowcolor{accent!8}
Prototype finding & Practical effect & Design response \\
\midrule
Sequential workflow & Runs became long and hard to inspect when many hosting
targets existed. & Use an orchestrator with explicit route state. \\
Weak browser continuity & Browser context could be lost between steps. & Use
profile-scoped browser sessions through tools. \\
Manual review & Failures required manual reconstruction from workflow output. &
Persist runtime events, screenshots, model calls, and tool calls. \\
No batch regression & Prompt changes were judged from small examples. & Build
dataset-backed batch evaluation. \\
\bottomrule
\end{tabularx}
\caption{n8n prototype lessons and implementation responses.}
\label{tab:n8n-agent-lessons}
\end{table}

\begin{figure}[!htbp]
\centering
\resizebox{0.88\textwidth}{!}{%
\begin{tikzpicture}[
  prob/.style={draw=gray!60, fill=gray!8, rounded corners=5pt,
    minimum width=3.75cm, minimum height=0.82cm, align=center, font=\footnotesize\bfseries},
  sol/.style={draw=green!55!black, fill=green!8, rounded corners=5pt,
    minimum width=3.75cm, minimum height=0.82cm, align=center, font=\footnotesize\bfseries},
  arr/.style={->, >=Stealth, thick, accent!80, shorten <=3pt, shorten >=3pt},
  heading/.style={font=\scriptsize\bfseries, text=gray!70}
]

\node[heading] at (0,3.88) {Prototype limitation};
\node[heading] at (6.4,3.88) {OWC design response};

\node[prob] (p1) at (0,3.0) {Sequential workflow};
\node[prob] (p2) at (0,1.8) {No durable memory};
\node[prob] (p3) at (0,0.6) {Fragile browser state};
\node[prob] (p4) at (0,-0.6) {Manual evaluation};

\node[sol] (s1) at (6.4,3.0) {Orchestrated specialists};
\node[sol] (s2) at (6.4,1.8) {Reusable memory hints};
\node[sol] (s3) at (6.4,0.6) {Profile-scoped MCP sessions};
\node[sol] (s4) at (6.4,-0.6) {Stored traces + regression routes};

\foreach \i in {1,2,3,4}
  \draw[arr] (p\i) -- (s\i);

\end{tikzpicture}
}
\caption{How n8n limitations shaped the Open Web Catcher implementation.}
\label{fig:n8n-to-reference}
\end{figure}


\section{Open Web Catcher Platform}

\subsection{Runtime Boundary}

OWC is the maintained reference/platform boundary described in this report. It
groups the FastAPI backend, Next.js operator console, PostgreSQL persistence,
browser tool services, specialist-agent runtime, provider/channel enrichment,
and reviewable outputs.

\subsection{Platform Modules}

\begin{table}[!htbp]
\centering
\begin{tabularx}{\textwidth}{Q{3.0cm} Y Y}
\toprule
\rowcolor{accent!8}
Module & Implementation role & Reason for boundary \\
\midrule
FastAPI backend & Starts jobs, exposes run detail, resolves configuration, and
persists trace data~\cite{fastapi}. & Keeps runtime control outside the UI. \\
Next.js console & Launches runs, follows progress, reviews evidence, and
changes settings~\cite{nextjs}. & Gives operators one review surface. \\
Browser tools & Inspect pages, click controls, capture screenshots, and
harvest network evidence~\cite{mcp-spec,puppeteer}. & Keeps browser actions
explicit and logged. \\
Agent runtime & Runs classification, landing, hosting, embedded, provider, and
email stages. & Splits responsibilities into bounded tasks. \\
PostgreSQL & Stores runs, events, streams, screenshots, provider rows, drafts,
model calls, and tool calls. & Makes results auditable. \\
\bottomrule
\end{tabularx}
\caption{Implementation modules of OWC.}
\label{tab:implementation-areas}
\end{table}

\subsection{Relationship Between n8n and OWC}

n8n validated the browser-agent behavior and served as the backup deployment
artifact for the crawler workflow. OWC formalized the maintainable architecture:
clearer agents, stronger observability, typed persistence, richer evaluation,
and a more structured operator console.

\subsection{Implementation Areas}

Implementation work was divided across backend APIs, UI review surfaces,
browser tools, prompt contracts, persistence, provider/channel enrichment,
evaluation, and queue-driven deployment direction.

\section{Backend Implementation}

\subsection{FastAPI Run Control}

The backend starts investigations, resolves runtime settings, registers active
traces, and exposes APIs for workflow launch, single-agent tests, run detail,
configuration, pricing, provider lookup, and dataset execution.

Run control is implemented as the entry point of the evidence workflow. When a
URL is submitted, the backend creates the run record, attaches runtime
configuration, selects the agent route policy, and returns an identifier that
the console can use for progress and review. This boundary matters because the
browser-agent work can be slow: the UI should not own long-running browser
state, and the final run should remain inspectable even if the browser worker
fails later.

\subsection{Agent Execution}

Agent execution is organized around a route owned by the orchestrator. The
orchestrator prepares handoffs and calls specialists for classification,
landing discovery, hosting interaction, embedded-player inspection, and
downstream enrichment.

Each specialist receives a narrower task than the complete investigation. The
classification agent decides the page role; the landing agent searches for
hosting candidates; the hosting agent interacts with servers, popups, and
players; the embedded agent focuses on iframe/player-only contexts; and the
provider/email stage packages downstream context. The backend records the
input, output, status, and duration of each stage so that a failed run can be
debugged at the responsible boundary.

\subsection{Trace Assembly}

Run detail is assembled from multiple records: status, events, screenshots,
streams, tool calls, model usage, provider rows, channel context, and generated
drafts. This is why the backend acts as a review API, not only a job launcher.

The run-detail response is intentionally composite. A reviewer needs one screen
that answers several questions at once: what page was tested, which route was
selected, which browser actions were attempted, what the browser showed, which
stream evidence was found, how provider lookup behaved, and what the run cost.
The backend therefore assembles normalized persistence rows into a coherent
evidence bundle instead of returning only the final agent text.

\subsection{Configuration Access}

Runtime configuration includes model choices, pricing data, browser settings,
tool availability, provider keys, memory behavior, and dataset settings.

Configuration is read by the backend before a run starts and again by review
surfaces that explain model and pricing behavior. Keeping this access in the
backend avoids hardcoded frontend assumptions about available models, tool
profiles, provider keys, or pricing rows. It also makes it possible to compare
two runs by the settings that were active when they executed.

\subsection{Error Handling}

Errors are preserved as runtime events and stage outputs. The system favors a
clear failed or partial state over a confident answer without evidence.

This error policy is important for evidence work. A failed harvest, blocked
page, popup trap, missing provider contact, or timeout is still useful if it is
stored with the responsible stage and visible context. The implementation
therefore avoids hiding failures behind a generic completed status. It returns
partial or failed outcomes with enough trace data for an engineer or reviewer
to decide whether a rerun, prompt fix, tool fix, or manual review is needed.

\section{Operator Console Implementation}

\subsection{Run Launch Page}

The launch surface accepts target URLs or dataset batches and starts
traceable investigations.

The launch page is the operator's controlled entry point. It separates single
URL testing from dataset-backed batch execution, because these two modes serve
different needs. A single URL is useful for debugging a page family or a
regression case. A dataset batch is useful for comparing prompt, model, or tool
changes against previously working targets. Both modes create the same kind of
run records so their outputs remain comparable.

\subsection{Run Detail Page}

The run detail page shows status, stage outputs, screenshots, streams,
provider/channel context, tool history, model usage, and errors.

The detail page is designed around trace reconstruction. A reviewer should be
able to start from the final status and work backward to the classification
decision, specialist handoff, browser actions, screenshots, stream rows, and
provider enrichment. This view is where a run becomes reviewable: it connects
the final result to the concrete evidence and failure path.

\subsection{Evidence Review Page}

The evidence review view groups screenshots, media candidates, source pages,
provider rows, and reviewable email drafts so a human can decide whether the
run is usable.

This page avoids treating every successful technical extraction as operationally
accepted. The operator still checks whether screenshots support the claim,
whether stream or embedded evidence belongs to the investigated page, whether
provider context is useful, and whether an email draft is complete enough for
human approval. The system prepares the package, but it does not replace the
reviewer.

\subsection{Telemetry Dashboard}

The overview dashboard gives a compact operational snapshot: runs, strict
stream success, stream yield, total cost, total tokens, cached input, LLM
calls, tool success, recent events, active runs, and persistence health.

\begin{figure}[!htbp]
\centering
\begingroup
\setlength{\fboxsep}{3pt}
\setlength{\fboxrule}{0.35pt}
\fcolorbox{gray!35}{white}{\includegraphics[width=0.96\textwidth]{dashboard-overview.png}}
\endgroup
\caption{Operator console overview with run, stream-yield, token, cost, tool, and persistence telemetry.}
\label{fig:dashboard-overview}
\end{figure}

\subsection{Provider and Channel Review View}

Provider and channel views make downstream context explicit: detected channel
labels, stream hosts, provider names, countries, contacts when available, and
missing-data flags.

The implementation keeps channel detection, provider resolution, and abuse
contact discovery separate. A run may detect a broadcaster label but lack a
provider contact, or it may resolve a provider without a clear channel. Showing
these fields separately prevents the review interface from implying certainty
that the evidence does not support.

\subsection{Runtime Configuration View}

Runtime settings expose model selection, provider configuration, tool health,
and cost/pricing context for engineers.

This view supports maintenance. Engineers can check which model is active,
which provider catalog or pricing rows are available, which browser tools are
registered, and whether cost telemetry is being captured. That context matters
when an evaluation run changes behavior after a model, prompt, or runtime
configuration update.

\begin{figure}[!htbp]
\centering
\begin{minipage}{0.48\textwidth}
  \centering
  \begingroup\setlength{\fboxsep}{2pt}\setlength{\fboxrule}{0.3pt}%
  \fcolorbox{gray!35}{white}{\includegraphics[width=0.94\linewidth]{frontend-dashboard.png}}\endgroup\\[-0.1em]
  {\scriptsize\bfseries Dashboard telemetry}
\end{minipage}\hfill
\begin{minipage}{0.48\textwidth}
  \centering
  \begingroup\setlength{\fboxsep}{2pt}\setlength{\fboxrule}{0.3pt}%
  \fcolorbox{gray!35}{white}{\includegraphics[width=0.94\linewidth]{frontend-live-workflow.png}}\endgroup\\[-0.1em]
  {\scriptsize\bfseries Live Pipeline launcher}
\end{minipage}

\vspace{0.55em}
\begin{minipage}{0.48\textwidth}
  \centering
  \begingroup\setlength{\fboxsep}{2pt}\setlength{\fboxrule}{0.3pt}%
  \fcolorbox{gray!35}{white}{\includegraphics[width=0.94\linewidth]{frontend-runs-websites.png}}\endgroup\\[-0.1em]
  {\scriptsize\bfseries View Results dataset}
\end{minipage}\hfill
\begin{minipage}{0.48\textwidth}
  \centering
  \begingroup\setlength{\fboxsep}{2pt}\setlength{\fboxrule}{0.3pt}%
  \fcolorbox{gray!35}{white}{\includegraphics[width=0.94\linewidth]{frontend-providers.png}}\endgroup\\[-0.1em]
  {\scriptsize\bfseries Provider Results}
\end{minipage}

\vspace{0.55em}
\begin{minipage}{0.48\textwidth}
  \centering
  \begingroup\setlength{\fboxsep}{2pt}\setlength{\fboxrule}{0.3pt}%
  \fcolorbox{gray!35}{white}{\includegraphics[width=0.94\linewidth]{frontend-dashboard-tools.png}}\endgroup\\[-0.1em]
  {\scriptsize\bfseries Tool telemetry}
\end{minipage}\hfill
\begin{minipage}{0.48\textwidth}
  \centering
  \begingroup\setlength{\fboxsep}{2pt}\setlength{\fboxrule}{0.3pt}%
  \fcolorbox{gray!35}{white}{\includegraphics[width=0.94\linewidth]{frontend-settings-models.png}}\endgroup\\[-0.1em]
  {\scriptsize\bfseries Model settings}
\end{minipage}

\caption{Current Next.js operator-console screens used for launch, review, telemetry, provider enrichment, and runtime configuration.}
\label{fig:frontend-console-screenshots}
\end{figure}


\section{Browser Tool Implementation}

\subsection{Navigation Tool Implementation}

Navigation tools open targets, follow direct navigation, return to prior pages,
scroll to reveal content, and wait for page state. They return final URL,
readiness signals, and screenshots where relevant.

\subsection{Inspection Tool Profiles}

Inspection tools are split between lightweight context reads, targeted element
queries, frame/media detail, classification-oriented inspection, and
profile-specific landing/hosting/embedded inspection.

\subsection{Interaction Tool Implementation}

Interaction tools perform clicks, typing, selection, checkbox/radio actions,
media play attempts, gestures, coordinates, and the higher-level
\texttt{interact} flow. Interaction is followed by verification rather than
accepted as proof.

\subsection{Screenshot Capture}

Screenshots are captured for classification, player activation, evidence,
popup, blocker, and failure states. They are stored with the run so reviewers
can inspect the browser-visible claim.

\subsection{Stream and Manifest Harvesting}

\texttt{capture\_streams} and \texttt{harvest} capture \texttt{m3u8},
\texttt{mpd}, \texttt{mp4}, \texttt{webm}, iframe, DOM, player-object,
performance-entry, and network evidence after playback or player readiness is
observed.

\subsection{Popup, Tab, and Frame Handling}

The browser runtime records tab changes, closes unrelated popups when needed,
restores focus, and uses frame paths so tools operate on the intended page
context.

\subsection{Browser Profile-Scoped Sessions}

Profile-scoped sessions preserve state for a run while avoiding uncontrolled
sharing across unrelated investigations.

\begin{figure}[!htbp]
\centering
\resizebox{0.88\textwidth}{!}{%
\begin{tikzpicture}[
  agent/.style={draw=orange!80!black, fill=orange!9, rounded corners=5pt,
    minimum width=3.0cm, minimum height=0.82cm, align=center, font=\footnotesize\bfseries},
  tool/.style={draw=purple!65!black, fill=purple!7, rounded corners=5pt,
    minimum width=3.1cm, minimum height=0.82cm, align=center, font=\footnotesize\bfseries},
  browser/.style={draw=gray!65, fill=gray!10, rounded corners=5pt,
    minimum width=3.1cm, minimum height=0.82cm, align=center, font=\footnotesize\bfseries},
  result/.style={draw=green!55!black, fill=green!8, rounded corners=5pt,
    minimum width=3.2cm, minimum height=0.82cm, align=center, font=\footnotesize\bfseries},
  arr/.style={->, >=Stealth, thick, accent!80, shorten <=3pt, shorten >=3pt},
  ret/.style={->, >=Stealth, thick, gray!65, dashed, shorten <=3pt, shorten >=3pt},
  lab/.style={font=\tiny\bfseries, fill=white, inner sep=1.2pt, text=gray!75!black}
]

\node[agent] (agent) at (0,2.1) {Specialist agent\\requests action};
\node[tool] (mcp) at (4.1,2.1) {MCP profile\\endpoint};
\node[browser] (chrome) at (8.2,2.1) {Puppeteer\\Chrome profile};
\node[result] (result) at (8.2,0.15) {Structured result\\state + evidence};
\node[result] (trace) at (4.1,0.15) {Trace update\\tool + memory};
\node[agent] (next) at (0,0.15) {Next agent turn\\verify or stop};

\draw[arr] (agent) -- node[lab, above]{tool call} (mcp);
\draw[arr] (mcp) -- node[lab, above]{browser command} (chrome);
\draw[ret] (chrome.south) -- node[lab, right]{observation} (result.north);
\draw[ret] (result) -- (trace);
\draw[ret] (trace) -- (next);
\draw[ret] (next.north) -- ++(0,0.65) -- (agent.south);

\end{tikzpicture}
}
\caption{Browser tool execution loop through MCP and Puppeteer.}
\label{fig:mcp-tool-loop}
\end{figure}


\section{Agents, Prompts, Memory, and Models}

\subsection{Prompt Compilation}

Each browser-facing agent compiles a role-specific prompt from the base policy,
agent contract, runtime settings, task URL, handoff context, and memory hints.

\begin{figure}[!htbp]
\centering
\resizebox{0.94\textwidth}{!}{%
\begin{tikzpicture}[
  box/.style={draw=accent, fill=blue!5, rounded corners=5pt,
    minimum width=3.8cm, minimum height=0.9cm, align=center, font=\footnotesize\bfseries},
  weak/.style={draw=orange!70, fill=orange!10, rounded corners=5pt,
    minimum width=3.8cm, minimum height=0.9cm, align=center, font=\footnotesize\bfseries},
  stop/.style={draw=gray!65, fill=gray!10, rounded corners=5pt,
    minimum width=3.8cm, minimum height=0.9cm, align=center, font=\footnotesize\bfseries},
  ok/.style={draw=green!55!black, fill=green!8, rounded corners=5pt,
    minimum width=3.8cm, minimum height=0.9cm, align=center, font=\footnotesize\bfseries},
  arr/.style={->, >=Stealth, thick, accent!80, shorten <=3pt, shorten >=3pt},
  loop/.style={->, >=Stealth, thick, gray!70, dashed, shorten <=3pt, shorten >=3pt},
  lab/.style={font=\tiny\bfseries, fill=white, inner sep=1.2pt, text=gray!70}
]
\node[weak] (old1) at (0,3.05) {Intent prompt};
\node[weak] (old2) at (4.2,3.05) {One-pass tool use};
\node[stop] (old3) at (8.4,3.05) {Premature answer};
\draw[arr] (old1)--(old2);
\draw[arr] (old2)--(old3);

\node[box] (new1) at (0,0.8) {Observe page state};
\node[box] (new2) at (4.2,0.8) {Reason + choose tool};
\node[box] (new3) at (8.4,0.8) {Act with MCP tool};
\node[ok] (new4) at (12.6,0.8) {Validate evidence};
\node[ok] (new5) at (12.6,-1.25) {Return result\\or clear failure};
\draw[arr] (new1)--(new2);
\draw[arr] (new2)--(new3);
\draw[arr] (new3)--(new4);
\draw[arr] (new4)--(new5);
\draw[loop] (new4.north) -- ++(0,0.78) -| node[lab, above]{not enough evidence} (new2.north);
\end{tikzpicture}
}
\caption{Prompt evolution from one-pass intent prompts to evidence-driven agent loops.}
\label{fig:prompt-evolution}
\end{figure}


\subsection{Agent Prompt Contracts}

\begin{table}[!htbp]
\centering
\begin{tabularx}{\textwidth}{Q{2.8cm} Y Y Y}
\toprule
\rowcolor{accent!8}
Agent & Prompt purpose & Required output & Main tool rules \\
\midrule
Orchestrator & Route the run and aggregate specialist results. &
\texttt{PipelineResult} with status, evidence, and failure reason. & Does not
browse directly; routes by page type and evidence state. \\
Classification & Decide whether the page is landing, hosting, embedded,
blocked, or other. & Page type, confidence, and route reason. & Inspect and
screenshot only; no interactions. \\
Landing & Discover events, channels, match cards, and hosting candidates. &
Ranked candidate list with context. & Use landing inspect, reveal actions, and
screenshots; avoid whole-domain crawling. \\
Hosting & Activate players, try servers, recover popups, and capture streams or
embedded handoffs. & Streams, embedded URL, screenshots, diagnostics, or clear
failure. & Inspect before acting; verify after acting; harvest after readiness. \\
Embedded & Inspect player-only or iframe contexts and recover delayed manifests.
& Streams, player type, screenshots, or blocked-player reason. & Stay in the
assigned player context and avoid unrelated navigation. \\
Provider/email & Enrich stream hosts and prepare reviewable drafts. & Provider
groups, channel context, draft email, missing-data flags. & Do not browse,
invent contacts, or send emails automatically. \\
\bottomrule
\end{tabularx}
\caption{Compact agent prompt contract summary.}
\label{tab:agent-prompt-contract-summary}
\end{table}

\subsection{Structured JSON Outputs}

Agent outputs are structured so the backend can persist them and the console
can render them consistently. Missing evidence is represented as an explicit
failure reason instead of omitted silently.

\subsection{Short Memory}

Memory is a hint layer. It can provide remembered selectors, pagination
patterns, navigation playbooks, server records, screenshots, stream URLs,
failure cues, and continuation notes, but the current browser state remains
the source of truth.

\subsection{Layout Hints}

Layout hints help agents avoid repeated broad inspection when a page family is
known. They are still validated against the live DOM and screenshot before use.

\subsection{Model Selection}

The selected model had to support long browser traces, reliable tool calls,
reasonable latency, and repeated regression runs at observable cost. Rate
limits, cached input, token usage, and pricing were therefore tracked during
evaluation~\cite{gemini-rate-limits,gemini-pricing,gemini-caching}.

\subsection{Token and Cost Tracking}

Each model call records model name, input tokens, cached input tokens, output
tokens, total tokens, and estimated spend when available.

\section{Persistence and Observability}

The persistence layer was implemented to answer two different review needs.
Evidence records support a human review decision: what was observed, what
stream or embedded-player evidence was collected, what screenshots support it,
which channel/provider context exists, and what draft was prepared. Telemetry
records support engineering diagnosis: which agent ran, which tool was called,
which model generated the decision, what failed, how long it took, and how much
it cost.

\begin{figure}[!htbp]
\centering
\resizebox{\textwidth}{!}{%
\begin{tikzpicture}[
  node distance=1.05cm,
  core/.style={draw=accent, fill=blue!6, rounded corners=4pt,
    minimum width=3.35cm, minimum height=1.0cm, align=center, font=\small\bfseries},
  evidence/.style={draw=green!45!black, fill=green!6, rounded corners=3pt,
    minimum width=3.25cm, minimum height=0.82cm, align=center, font=\footnotesize},
  telemetry/.style={draw=orange!70!black, fill=orange!8, rounded corners=3pt,
    minimum width=3.25cm, minimum height=0.82cm, align=center, font=\footnotesize},
  review/.style={draw=gray!70, fill=gray!8, rounded corners=4pt,
    minimum width=3.35cm, minimum height=1.0cm, align=center, font=\small\bfseries},
  arrow/.style={->, >=Stealth, thick, accent}
]

\node[core] (run) {Run identity\\target, status, route};

\node[evidence, right=1.45cm of run, yshift=1.95cm] (screens) {Screenshots\\evidence and failure};
\node[evidence, below=0.28cm of screens] (streams) {Streams and embeds\\m3u8, mpd, mp4, iframe};
\node[evidence, below=0.28cm of streams] (provider) {Provider and channel\\host, IP, contact status};
\node[evidence, below=0.28cm of provider] (drafts) {Review drafts\\not sent automatically};

\node[telemetry, right=1.35cm of screens] (events) {Runtime events\\ordered timeline};
\node[telemetry, below=0.28cm of events] (tools) {Tool calls\\action, status, duration};
\node[telemetry, below=0.28cm of tools] (models) {Model calls\\prompt, tokens, cost};
\node[telemetry, below=0.28cm of models] (errors) {Errors and memory\\failure cues, hints};

\node[review, right=1.45cm of tools] (console) {Operator console\\run detail and evidence};
\node[review, above=0.9cm of console] (dashboard) {Telemetry dashboard\\cost, tools, success};
\node[review, below=0.9cm of console] (regression) {Regression analysis\\compare old and new runs};

\draw[arrow] (run) -- (screens);
\draw[arrow] (run) -- (streams);
\draw[arrow] (run) -- (provider);
\draw[arrow] (run) -- (drafts);
\draw[arrow] (screens) -- (events);
\draw[arrow] (streams) -- (tools);
\draw[arrow] (provider) -- (models);
\draw[arrow] (drafts) -- (errors);
\draw[arrow] (events) -- (dashboard);
\draw[arrow] (tools) -- (console);
\draw[arrow] (models) -- (console);
\draw[arrow] (errors) -- (regression);
\draw[arrow] (console) -- (dashboard);
\draw[arrow] (console) -- (regression);

\node[font=\footnotesize, align=center, text=gray!70]
  at ($(events.north)!0.5!(models.south)+(0.2,0)$)
  {};

\end{tikzpicture}%
}
\caption{Persistence and observability flow: evidence records support review, while telemetry records support diagnosis and regression.}
\label{fig:persistence-review-flow}
\end{figure}


\subsection{Stored Evidence Records}

Stored evidence records are the reviewable artifacts of the investigation.
They include screenshots, embedded URLs, stream or manifest candidates,
channel labels, provider rows, and generated email drafts. Each record keeps a
connection to the run and, where possible, to the page URL, frame context,
timestamp, and tool that produced it. This prevents the final evidence bundle
from becoming a loose list of URLs with no explanation.

The most important implementation choice is that evidence is stored even when
it is partial. A blocked screenshot, a failed player state, or a provider row
without an abuse contact can still explain why a run did not reach strict
success. For review, a traceable failure is better than a missing row.

\subsection{Stored Telemetry Records}

Telemetry records are the engineering record of the run. They include agent
runs, runtime events, model calls, tool calls, prompt compilations, errors,
and memory hints. These rows are not shown to the operator as raw logs only;
they are summarized into run detail, dashboard metrics, and regression views.

Separating telemetry from evidence keeps the review model clean. A reviewer
needs screenshots and stream context first, while an engineer needs event
order, tool status, model cost, and error detail. The same run therefore feeds
two interfaces: evidence review and system diagnosis.

\subsection{Runtime Event Logging}

Runtime events preserve the investigation timeline and make failed runs
diagnosable. Events are small and timestamped: run started, agent started,
tool requested, tool completed, screenshot captured, stream found, provider
resolved, failure recorded, and run completed. This timeline explains the
difference between superficially similar failures. For example,
\texttt{no\_streams} can mean that the player never activated, the harvest ran
too early, the browser was on a popup tab, or the target was blocked.

\begin{table}[!htbp]
\centering
\begin{tabularx}{\textwidth}{Q{3.6cm}Y}
\toprule
\rowcolor{accent!8}
Event type & Review or debugging value \\
\midrule
Run started & Target URL, selected model, browser profile, and batch source. \\
Agent started & Specialist name, handoff reason, and evidence objective. \\
Tool requested & Tool name, input summary, active page URL, and current stage. \\
Tool completed & Status, duration, output summary, and error if present. \\
Screenshot captured & Screenshot path, page URL, and capture reason. \\
Stream found & Stream host, source tool, source page, and frame context. \\
Provider resolved & Host, IP, provider name, country, and lookup source. \\
Failure recorded & Responsible stage, visible blocker, and failure reason. \\
Run completed & Final status, evidence counts, cost, and acceptance state. \\
\bottomrule
\end{tabularx}
\caption{Runtime event types used to reconstruct an investigation run.}
\label{tab:runtime-event-types-main}
\end{table}

\subsection{Tool Call Logging}

Tool calls store the browser action or observation, status, duration, active
URL, input summary, and output summary. This makes browser behavior auditable.
When a click opens a popup, the trace can show the click, the new context, the
focus restoration, and the following verification step. When harvesting
returns no stream, the tool-call order can show whether media readiness was
checked first.

\begin{table}[!htbp]
\centering
\begin{tabularx}{\textwidth}{Q{3.4cm}Y Y}
\toprule
\rowcolor{accent!8}
Tool family & What the implementation stores & Failure pattern it exposes \\
\midrule
Inspection & Page profile, candidates, frame paths, media state, and output
size. & Broad context repeated without state change, or wrong page profile. \\
Interaction & Locator strategy, click/play/type target, status, and screenshot
checkpoint. & Popup focus loss, wrong control, or unverified action. \\
Navigation & Intended target, final URL, redirects/status, and readiness
result. & Accidental reload, lost player state, or challenge page. \\
Screenshot & Capture type, page URL, reason, and stored path. & Evidence
claim without a visual checkpoint. \\
Harvest & Detection layer, stream rows, network diagnostics, iframe
diagnostics, and video state. & Harvest before playback or missing frame
context. \\
Provider lookup & Host, resolved IP, provider name, country, contact status,
and lookup source. & Provider found but abuse contact missing. \\
\bottomrule
\end{tabularx}
\caption{Tool-call logging fields and the failures they make visible.}
\label{tab:tool-call-review-main}
\end{table}

\subsection{Model Call Logging}

Model calls store prompt metadata, model selection, tokens, cached input, and
cost. During evaluation, this made it possible to compare prompt versions by
cost per useful run rather than by total spend alone. The implementation keeps
model provider, model name, prompt version, input tokens, cached input tokens,
output tokens, and estimated cost when those values are available. This is
also how the dashboard can explain why a regression run became more expensive
after a prompt or tool-policy change.

\subsection{Screenshot Storage}

Screenshots are stored as evidence and as debugging context for failures. The
implementation uses screenshots at meaningful state changes: classification,
landing candidate discovery, hosting activation, embedded-player inspection,
blocked pages, popup states, and final evidence or failure moments. The goal is
not to capture every wait loop; the goal is to preserve the visual state that
supports a claim or explains a failure.

\subsection{Error and Failure Recording}

Failures remain visible as structured status, stage, and reason fields. A run
can be \texttt{success}, \texttt{partial}, \texttt{failed},
\texttt{no\_streams}, \texttt{no\_hosting\_pages}, or \texttt{blocked}, but
the status alone is not enough. The implementation also stores the responsible
stage, visible blocker, last useful screenshot, last tool call, and error
message. This makes a failure actionable: the next fix can target the
classifier, landing discovery, hosting interaction, embedded harvest,
provider lookup, or runtime configuration.

\begin{table}[!htbp]
\centering
\begin{tabularx}{\textwidth}{Q{3.2cm} Y Y}
\toprule
\rowcolor{accent!8}
Stored record & Example question answered & Why it matters \\
\midrule
Runtime event & Which stage failed and why? & Supports failure diagnosis. \\
Screenshot & What did the browser actually show? & Supports visual review. \\
Tool call & Which browser action was attempted? & Separates prompt errors from
tool/runtime errors. \\
Model call & Which model was used and how many tokens were spent? & Supports
cost and model evaluation. \\
Provider/channel row & Which host, provider, or channel context was detected? &
Supports rights-enforcement review. \\
\bottomrule
\end{tabularx}
\caption{Observability records used during run review.}
\label{tab:persistence-observability}
\end{table}

\section{Provider, Channel Detection, and Email Preparation}

\subsection{Channel and Broadcaster Detection}

The workflow attempts to detect broadcaster/channel indicators from visible
page labels, DOM text, screenshots, stream context, and metadata. This keeps
the evidence model broader than a single rights holder while still supporting
the beIN SPORTS context when beIN indicators are detected.

\subsection{Stream Host Extraction}

Stream URLs and embedded-player URLs are normalized to hosts when possible so
provider enrichment can resolve infrastructure context.

\subsection{Provider Lookup}

Provider lookup uses host/IP and registration sources such as IP and RDAP data
to enrich stream evidence~\cite{ipinfo-api,icann-rdap}.

\subsection{Rights-Holder Email Context}

When channel context can be identified, the email stage can prepare a draft
using the corresponding rights-holder or contact context. When it cannot, the
missing channel context remains explicit.

\subsection{Reviewable Email Drafts}

Generated messages are drafts only. They group evidence, screenshots,
provider/channel context, and run identifiers for human review.

\subsection{Human Review Boundary}

The system does not send takedown emails automatically. A human reviewer
validates the evidence and decides downstream use.

\section{Deployment Direction}

\subsection{Company Crawler Environment}

The existing company crawler environment remains the operational baseline for
known patterns. The browser-agent work supports this environment instead of
claiming to replace it.

\subsection{n8n Landing Agent as Backup Solution}

The n8n landing agent was deployed as a backup to the existing crawler
workflow. When the crawler failed to recover useful landing-page candidates
after site changes, the n8n agent could inspect the page more flexibly and
return candidate URLs or evidence metadata for review or continuation.

\subsection{\texttt{agent\_tasks} Input Queue}

The deployment direction uses an Azure Service Bus input queue named
\texttt{agent\_tasks} for target URLs and task metadata.

\subsection{Browser Execution Container}

The browser execution container runs the workflow and Puppeteer tooling close
to the browser session so long investigations can run asynchronously.

\subsection{\texttt{agent\_results} Output Queue}

Results are returned through \texttt{agent\_results}: matched URLs, candidate
metadata, evidence references, and status. Queueing fits slow browser runs
because work can be retried and consumed asynchronously~\cite{azure-servicebus}.

\subsection{Relationship Between Backup Deployment and OWC}

The backup deployment is operationally useful for the crawler context. OWC is
the maintainable architecture that generalizes the approach with stricter
agents, richer observability, typed persistence, and repeatable evaluation.

\begin{figure}[!htbp]
\centering
\includegraphics[width=0.82\textwidth,height=0.58\textheight,keepaspectratio,trim=8 8 8 8,clip]{figures/deployment-physical-architecture.png}
\caption[Company deployment flow for the landing-agent backup.]{Company deployment flow for the landing-agent backup, from the primary workflow and Azure Service Bus queues to the focused Azure Container Apps job, browser runtime, external dependencies, and result payload.}\label{fig:physical-architecture}
\end{figure}


\section{Implementation Summary}

The implementation moved from a known crawler baseline and n8n validation
toward a reviewable browser-agent architecture. The main engineering gain was
not a single extraction script. It was a runtime where prompt changes, browser
tools, screenshots, costs, stream evidence, channel/provider data, and failures
can be inspected together.
